% !TEX root = ../main.tex
%----------------------------------------------------------
% 07-worked-examples.tex — Section 6: Worked Examples
% NEW content (figure-heavy) — structure with placeholders
%----------------------------------------------------------

\section{Worked Examples}\label{sec:worked_examples}

The preceding sections developed three verification lanes---the native
scorecard, the possibility-to-probability bridge with information gain,
and the interval log score---in the abstract.  This section applies all
three to a concrete trio of stylized convective outlook scenarios.  Each
scenario is designed to isolate a different diagnostic situation: a sharp
forecast that verifies correctly, a hedged forecast that also verifies
correctly, and a sharp forecast that misses badly.  Walking through the
full scorecard, gauge visualization, bridge computation, and ILS
calculation for each case demonstrates how the metrics complement one
another and how subnormality propagates through every stage of the
framework.

%----------------------------------------------------------
\subsection{Scenario Setup}\label{sec:scenario_setup}

Three stylized possibilistic convective outlooks are defined, each
representing a distinct forecasting situation.  All distributions are
defined over $\Omega = \{\spc{NONE},\, \spc{MRGL},\, \spc{SLGT},\,
\spc{ENH},\, \spc{MDT},\, \spc{HIGH}\}$ with $K = 6$.  We consider a
notional point forecast extracted at Norman, Oklahoma: the SPC categorical
outlook at that location on each day serves as the observation.

\paragraph{Scenario A: Sharp-correct \spc{MDT} forecast.}

A textbook severe weather environment is in place: surface-based CAPE
exceeds 3000~J~kg$^{-1}$, 0--6~km bulk shear surpasses 40~kt, and a
well-defined dryline provides a reliable mesoscale trigger.  The
ingredients for significant severe thunderstorms are all present and
co-located, so the system assigns high possibility to \spc{MDT} while
leaving only token values for neighbouring categories.  Subnormality is
low ($m = 0.90$, $\ign = 0.10$): the system is nearly fully committed.

\[
\pi_A(\omega) = \begin{cases}
0.00 & \omega = \spc{NONE} \\
0.00 & \omega = \spc{MRGL} \\
0.05 & \omega = \spc{SLGT} \\
0.15 & \omega = \spc{ENH}  \\
0.90 & \omega = \spc{MDT}  \\
0.10 & \omega = \spc{HIGH}
\end{cases}
\]
Observed outcome: $c_{\mathrm{obs}} = \spc{MDT}$.

\paragraph{Scenario B: Hedged-correct \spc{ENH} forecast.}

The environment contains competing signals: moderate CAPE
(1500--2000~J~kg$^{-1}$) and adequate deep-layer shear, but the timing
and location of the mesoscale trigger remain uncertain.  Whether
convection organizes into discrete supercells or a mesoscale convective
system will determine whether the day verifies at \spc{SLGT},
\spc{ENH}, or even \spc{MDT}.  Faced with this ambiguity, the system
hedges: possibility is spread across four categories, and subnormality is
high ($m = 0.55$, $\ign = 0.45$), explicitly advertising the system's
limited confidence.

\[
\pi_B(\omega) = \begin{cases}
0.10 & \omega = \spc{NONE} \\
0.10 & \omega = \spc{MRGL} \\
0.40 & \omega = \spc{SLGT} \\
0.55 & \omega = \spc{ENH}  \\
0.30 & \omega = \spc{MDT}  \\
0.00 & \omega = \spc{HIGH}
\end{cases}
\]
Observed outcome: $c_{\mathrm{obs}} = \spc{ENH}$.

\paragraph{Scenario C: Sharp-wrong \spc{NONE} forecast.}

On the surface this appears a quiescent day: large-scale forcing is weak,
surface-based CAPE is marginal (less than 1000~J~kg$^{-1}$), and model
guidance unanimously favours no severe weather at the point of interest.
The system confidently assigns high possibility to \spc{NONE}---the true
null---and near-zero to all severe categories.  In reality, an unresolved
mesoscale convergence boundary interacts with the residual instability,
triggering an unexpected severe outbreak that verifies at \spc{MDT}---a
category the system deemed essentially impossible.

\[
\pi_C(\omega) = \begin{cases}
0.85 & \omega = \spc{NONE} \\
0.10 & \omega = \spc{MRGL} \\
0.05 & \omega = \spc{SLGT} \\
0.00 & \omega = \spc{ENH}  \\
0.00 & \omega = \spc{MDT}  \\
0.00 & \omega = \spc{HIGH}
\end{cases}
\]
Observed outcome: $c_{\mathrm{obs}} = \spc{MDT}$.

%----------------------------------------------------------
\subsection{Filling the Gauge}\label{sec:filling_gauge}

The gauge visualization maps each SPC category to a horizontal bar
whose fill level represents raw possibility $\pi(\omega)$.  A fully
filled bar ($\pi = 1.0$) would indicate complete plausibility; the gap
between the tallest bar and $1.0$ is the ignorance $\ign$, rendered as
empty space to the right of every bar.  The observed category is
highlighted in green so that the reader can immediately judge whether
the system placed high possibility on the outcome that actually
occurred.  This representation makes subnormality and its consequences
visible at a glance, without requiring the reader to parse numerical
tables.

Figure~\ref{fig:filling_gauge} shows the gauge visualization for
Scenario~A.  The horizontal bars are filled to the raw possibility
level; the gap from $\Pi_{\max}$ to $1.0$ visually encodes the
ignorance $\ign$.  The observed category (\spc{MDT}) is highlighted
in green.

\begin{figure}[t]
\centering
\includegraphics[width=\columnwidth]{fig3_filling_gauge.png}
\caption{Gauge visualization for Scenario~A (sharp-correct \spc{MDT}).
    Each horizontal bar represents an SPC category, filled to its raw
    possibility value $\pi(\omega)$.  The vertical dashed line marks
    $\Pi_{\max} = 0.90$; the gap to $1.0$ is ignorance $\ign = 0.10$.
    The observed category (\spc{MDT}, green) dominates the distribution,
    confirming high conditional necessity.}
\label{fig:filling_gauge}
\end{figure}

Were we to draw the same gauge for Scenario~B, the picture would change
qualitatively.  Five bars would be partially filled, with \spc{ENH} the
tallest at only $0.55$; the remaining bars for \spc{SLGT} ($0.40$) and
\spc{MDT} ($0.30$) would reach nearly the same height.  The large
$\ign = 0.45$ gap would dominate the right-hand side of every bar,
immediately signaling that the system is hedging rather than committing.
The observed category (\spc{ENH}, green) would still be the tallest bar,
but only marginally so---consistent with a forecast that is directionally
correct yet openly uncertain.

Scenario~C would present the most visually striking gauge.  A single bar
(\spc{NONE}) would dominate at $0.85$, with \spc{MRGL} contributing a
sliver at $0.10$ and the remaining four categories nearly empty.
The ignorance gap ($\ign = 0.15$) would be small, reinforcing the
impression of a confident forecast.  Yet the observed category
(\spc{MDT}, green) would appear as an unfilled bar at zero---a stark
visual indictment.  Placing the three gauges side by side reveals at a
glance that sharpness alone is not a virtue: it must be aimed at the
correct outcome.

%----------------------------------------------------------
\subsection{Full Scorecard Comparison}\label{sec:scorecard_comparison}

Table~\ref{tab:scorecard_comparison} collects the five-number scorecard
for all three scenarios.  Figure~\ref{fig:three_scenario} provides a
visual comparison, with scorecard values printed below each panel.

\begin{table}[htbp]
\centering
\caption{Five-number scorecard for the three worked scenarios.}
\label{tab:scorecard_comparison}
\small
\begin{tabular}{lccc}
\toprule
Statistic & Scenario A & Scenario B & Scenario C \\
    & (sharp-correct) & (hedged-correct) & (sharp-wrong) \\
\midrule
$m$ (subnormality)
    & $0.90$ & $0.55$ & $0.85$ \\
$\ign$ (ignorance)
    & $0.10$ & $0.45$ & $0.15$ \\
$\alpha^*$ (depth-of-truth)
    & $1.00$ & $1.00$ & $0.00$ \\
$\eta$ (nonspecificity)
    & $0.222$ & $0.439$ & $0.196$ \\
$\delta$ (resolution gap)
    & $+0.778$ & $+0.561$ & $-0.196$ \\
$\Nc^*$ (cond.\ necessity of truth)
    & $0.833$ & $0.273$ & $0.00$ \\
\bottomrule
\end{tabular}
\end{table}

\begin{figure*}[t]
\centering
\includegraphics[width=0.90\textwidth]{fig2_three_scenario.png}
\caption{Three-scenario comparison.  Each panel shows the raw possibility
    distribution $\pi(\omega)$ as bar charts; the observed category is
    highlighted in green, others in purple.  The five-number scorecard
    values are printed below each panel.  \textbf{(a)}~Scenario~A:
    sharp-correct \spc{MDT}---high $\alpha^*$, high $\Nc^*$, positive
    $\delta$.  \textbf{(b)}~Scenario~B: hedged-correct \spc{ENH}---high
    $\alpha^*$ but diffuse shape ($\eta$) and low $\Nc^*$.
    \textbf{(c)}~Scenario~C: sharp-wrong
    \spc{NONE}---$\alpha^* = 0$ (truth assigned zero possibility) and
    negative $\delta$.}
\label{fig:three_scenario}
\end{figure*}

Scenario~A represents the ideal case.  The depth-of-truth
$\alpha^* = 1.00$ confirms that the observed category (\spc{MDT})
received the highest raw possibility after normalization, and the
conditional necessity of truth $\Nc^* = 0.833$ indicates that
alternatives were strongly down-weighted.  Ignorance is low
($\ign = 0.10$), nonspecificity is modest ($\eta = 0.222$), and the
resolution gap $\delta = +0.778$ is strongly positive---the system knew
what was happening and said so confidently.  Every metric in the
scorecard rewards this forecast.

Scenario~B tells a different but equally honest story.  The
depth-of-truth remains perfect ($\alpha^* = 1.00$) because \spc{ENH},
the observed category, is indeed the possibilistic peak after shape
normalization.  However, the system's hedging shows up clearly:
$\ign = 0.45$ is nearly five times that of Scenario~A, nonspecificity
rises to $\eta = 0.439$, and $\Nc^* = 0.273$ is low because
\spc{SLGT} and \spc{MDT} were assigned plausible alternatives.  The
resolution gap $\delta = +0.561$ is positive but moderate, reflecting
genuine skill tempered by acknowledged uncertainty.  This is the
signature of a well-calibrated system on a difficult day.

Scenario~C is the pathological case.  The system assigned
$\pi(\spc{MDT}) = 0.00$, declaring the observed outcome impossible;
accordingly $\alpha^* = 0.00$ and $\Nc^* = 0.00$.  The resolution gap
$\delta = -0.196$ is the only negative value in the table, indicating
that the forecast pointed \emph{away} from the truth.  Critically,
ignorance was low ($\ign = 0.15$): the system was not hedging
cautiously---it was confident and wrong.  Low ignorance amplifies the
penalty because there is no uncertainty buffer to absorb the miss.

Contrasting Scenarios~B and~C reveals the scorecard's diagnostic power.
Both forecasts have noteworthy ignorance levels ($0.45$ and $0.15$,
respectively), yet the scorecard sharply distinguishes them.  Scenario~B
earns $\delta = +0.561$ because its hedging was aimed in the right
direction; Scenario~C earns $\delta = -0.196$ because its confidence
was aimed in the wrong direction.  A single summary statistic (e.g.,
mean possibility or maximum value) would conflate these two situations.
The five-number scorecard correctly separates ``honestly uncertain'' from
``confidently wrong,'' which is precisely the distinction that
operational forecasters need when tuning a possibilistic system.

%----------------------------------------------------------
\subsection{Pignistic Bridge Walkthrough}\label{sec:bridge_walkthrough}

We now apply the possibility-to-probability bridge
(Section~\ref{sec:poss_prob_bridge}) to each scenario, then compute
information gain relative to a uniform climatological baseline
$p_{\mathrm{clim}} = (1/6,\ldots,1/6) \approx 0.167$ per category over
the six SPC categories.  Recall the bridge steps: reserve ignorance mass
$p_{\mathrm{ign}} = \ign = 1 - m$; distribute the remaining $1 - \ign$
proportionally across outcomes via $p_i = \pi_i (1 - \ign) / \sum_j
\pi_j$; append the ignorance outcome.

\begin{itemize}[leftmargin=2em]
\item \textbf{Scenario~A:} $\ign = 0.10$, $\sum \pi = 1.20$.  The
    bridge yields $\mathbf{p}_A = (0.000,\, 0.000,\, 0.038,\, 0.113,\,
    0.675,\, 0.075,\, \underbrace{0.100}_{\text{ign}})$.  Since \spc{MDT}
    is observed, $-\!\log_2(0.675) = 0.567$~bits.  Climatological surprise
    is $-\!\log_2(1/6) = 2.585$~bits, giving
    $\mathrm{IG}_A = 2.585 - 0.567 = +2.018$~bits.
\item \textbf{Scenario~B:} $\ign = 0.45$, $\sum \pi = 1.45$.  The
    bridge yields $\mathbf{p}_B = (0.038,\, 0.038,\, 0.152,\, 0.209,\,
    0.114,\, 0.000,\, \underbrace{0.450}_{\text{ign}})$.  Since \spc{ENH}
    is observed, $-\!\log_2(0.209) = 2.258$~bits, giving
    $\mathrm{IG}_B = 2.585 - 2.258 = +0.327$~bits.
\item \textbf{Scenario~C:} $\ign = 0.15$, $\sum \pi = 1.00$.  The
    bridge yields $\mathbf{p}_C = (0.722,\, 0.085,\, 0.042,\, 0.000,\,
    0.000,\, 0.000,\, \underbrace{0.150}_{\text{ign}})$.  Since \spc{MDT}
    is observed, $p(\spc{MDT}) = 0.00$; using $\varepsilon = 0.01$ to
    regularize, $-\!\log_2(0.01) = 6.644$~bits, giving
    $\mathrm{IG}_C = 2.585 - 6.644 = -4.059$~bits.
\end{itemize}

The three IG values span a wide range and align with the scorecard
diagnostics.  Scenario~A's $+2.018$~bits reflects genuine skill: the
bridge concentrated probability on the correct outcome, and low ignorance
meant little mass was diverted to the ignorance bin.  Scenario~B's
$+0.327$~bits is positive but modest---the high $\ign = 0.45$ correctly
deflated $p(\spc{ENH})$ from what simple normalization would have
produced ($0.55 / 1.45 = 0.379$) down to $0.209$, honestly reflecting
the system's uncertainty.  Scenario~C incurs a catastrophic
$-4.059$~bits because the bridge assigned zero probability to the
observed outcome; even with the $\varepsilon$ floor, the logarithmic
penalty is severe.  Crucially, the low ignorance in Scenario~C
($\ign = 0.15$) provides almost no cushion: the system left itself no
escape route for an outcome it had ruled out.

Finally, we compute the interval log score (ILS) for the threshold
event $A_T = \{\spc{ENH},\, \spc{MDT},\, \spc{HIGH}\}$ across all three
scenarios.  For each forecast, the upper bound is $U = \max_{\omega \in
A_T} \pi(\omega)$; the shape-normalized distribution $\pinorm = \pi / m$
is used to compute conditional necessity via $\Nc(A_T) = 1 -
\max_{\omega \notin A_T} \pinorm(\omega)$; and the lower bound is $L = m
\cdot \Nc(A_T)$.  When the event occurs, the ILS penalty is
$-\!\log_2(U)$.

\begin{itemize}[leftmargin=2em]
\item \textbf{Scenario~A:} $U = \max(0.15, 0.90, 0.10) = 0.90$.
    Shape-normalizing gives $\pinorm = (0.000,\, 0.000,\, 0.056,\,
    0.167,\, 1.000,\, 0.111)$; $\max_{\omega \notin A_T} \pinorm =
    \max(0.000,\, 0.000,\, 0.056) = 0.056$, so $\Nc(A_T) = 0.944$ and
    $L = 0.90 \times 0.944 = 0.850$.  Event occurs (\spc{MDT}
    $\in A_T$): $\mathrm{ILS}_A = -\!\log_2(0.90) = 0.152$~bits.
\item \textbf{Scenario~B:} $U = \max(0.55, 0.30, 0.00) = 0.55$.
    Shape-normalizing gives $\pinorm = (0.182,\, 0.182,\, 0.727,\,
    1.000,\, 0.545,\, 0.000)$; $\max_{\omega \notin A_T} \pinorm =
    \max(0.182,\, 0.182,\, 0.727) = 0.727$, so $\Nc(A_T) = 0.273$ and
    $L = 0.55 \times 0.273 = 0.150$.  Event occurs (\spc{ENH}
    $\in A_T$): $\mathrm{ILS}_B = -\!\log_2(0.55) = 0.862$~bits.
\item \textbf{Scenario~C:} $U = \max(0.00, 0.00, 0.00) = 0.00$.
    Using $\varepsilon = 0.01$, the event occurs (\spc{MDT} $\in A_T$):
    $\mathrm{ILS}_C = -\!\log_2(0.01) = 6.644$~bits---maximum penalty.
\end{itemize}

The ILS results mirror and reinforce the bridge-based IG analysis.
Scenario~A earns the lightest penalty ($0.152$~bits) because the
system placed high possibility on the threshold event with tight
$[L, U] = [0.850,\, 0.900]$ bounds.  Scenario~B receives a moderate
penalty ($0.862$~bits): the system considered the event plausible
($U = 0.55$) but hedged substantially, yielding wide bounds
$[0.150,\, 0.550]$ that honestly reflect its ambivalence.  Scenario~C
incurs the maximum penalty ($6.644$~bits) because the system assigned
zero raw possibility to every category in $A_T$---the event was
declared impossible.  Taken together, the IG and ILS analyses confirm
that the scorecard diagnostics are not merely descriptive: they
translate directly into probabilistic and possibilistic skill
penalties that would drive appropriate system improvement.
